\documentclass[11pt,a4paper]{article}
\usepackage[utf8]{inputenc}
\usepackage[margin=1.0in]{geometry}
\usepackage{amsmath,amssymb,amsfonts,amsthm}
\usepackage{graphicx}
\usepackage{booktabs}
\usepackage{cite}
\usepackage{hyperref}
\usepackage{xcolor}
\usepackage{setspace}
\usepackage{abstract}

\hypersetup{
    colorlinks=true,
    linkcolor=navyblue,
    citecolor=navyblue,
    urlcolor=navyblue
}

\definecolor{navyblue}{rgb}{0.0, 0.15, 0.45}

\doublespacing

\title{\LARGE \textbf{Informational Memory Depth, Peer Contagion, and Spontaneous Phase Transitions in Emerging Financial Markets: A Non-Equilibrium Agent-Based Study of the Nifty 50}}

\author{
    \textbf{Anuj Kothari}\thanks{Department of Computer Science \& Finance Research Group, Indian Institute of Technology. Email: \texttt{anujkothari@example.com}.}
}

\date{\today}

\begin{document}

\maketitle

\begin{abstract}
\noindent We investigate the macroscopic emergence of financial herding, consensus lock-in, and market phase transitions resulting from micro-level cognitive memory horizons and peer sentiment contagion. Adapting non-equilibrium statistical mechanics to financial dynamics, we construct an agent-based framework where boundedly rational market participants ($N=150$) process historical daily return signals via finite memory queues of length $M$ subject to cognitive misinterpretation noise $\eta$. By calibrating the model against daily return series of the Indian Nifty 50 Index (2015--2025), we empirically reconstruct the effective potential landscape $V(n_+)$ governing the crowd order parameter $n_+(t)$. Our analytical and empirical results yield three fundamental conclusions. First, independent signal processing without social peer feedback ($p_{\text{peer}}=0$) is strictly restoring ($k>0$), establishing that market signals alone cannot induce herding due to Law-of-Large-Numbers suppression. Second, peer sentiment transmission ($p_{\text{peer}}=1.0$) triggers a supercritical pitchfork bifurcation at a critical cognitive memory depth $M_c \approx 9.82$ trading days, separating a unimodal, mean-reverting regime from a bistable, lock-in regime. Third, relaxation experiments starting from unanimous panic consensus ($n_+(0)=0$) reveal severe critical slowing down for conformist crowds as $M \to M_c$, whereas recency-weighted memory kernels ($b>1.0$) act as cognitive circuit breakers that cap relaxation times. Finally, we demonstrate that cognitive misinterpretation noise $\eta > 0.50$ inverts the sentiment feedback loop, converting the crowd into a self-stabilizing negative-feedback engine that inherently quenches speculative bubbles.

\vspace{0.3cm}
\noindent \textbf{Keywords:} Agent-Based Computational Economics, Behavioral Finance, Informational Memory, Pitchfork Bifurcation, Critical Slowing Down, Market Herding, Nifty 50.
\end{abstract}

\newpage

\section{Introduction}

The Classical Efficient Market Hypothesis (EMH) posits that asset prices instantaneously reflect all available public information, portraying price formation as a martingalian random walk driven strictly by exogenous fundamental shocks \cite{fama1970efficient}. However, decadal empirical evidence across global financial markets demonstrates ubiquitous non-linear anomalies, including heavy-tailed return distributions, volatility clustering, spontaneous speculative bubbles, and systemic market panics \cite{shiller2000irrational, lux1999scaling, cont2001empirical}. These macroscopic phenomena suggest that financial markets operate as complex, non-equilibrium open systems driven by endogenous social interactions and boundedly rational cognitive processing \cite{bouchaud2013more, farmer2009economy}.

A primary catalyst of market instability is *herding behavior*---the propensity of individual market participants to suppress private information and align sentiment with aggregate crowd opinions \cite{kirman1993ants, bikhchandani1992theory, Avery1998}. While classic behavioral literature attributes herding to psychological biases such as availability heuristics and fear-of-missing-out (FOMO) \cite{kahneman1979prospect}, the precise mathematical conditions governing the transition from efficient mean-reverting price discovery to locked-in herding regimes remain an active domain of inquiry.

Recent advances in non-equilibrium physics by Li et al. \cite{li2026informational} established that autonomous swarm units equipped with finite informational memory queues undergo sharp supercritical pitchfork bifurcations. Above a critical memory threshold $M_c$, interaction dynamics alone cause a symmetric system to spontaneously break symmetry, locking the population into persistent macroscopic consensus states.

In this study, we bridge non-equilibrium statistical mechanics and behavioral market microstructure by constructing a memory-kernel agent-based model parameterized by high-frequency return series of the Indian Nifty 50 Index. We model a heterogeneous market crowd ($N=150$) characterized by five empirical trader archetypes: *Herding Conformists* (Pushovers), *Recency-Biased Traders* (Opportunists), *Anchored History Traders* (Traditionalists), *Mean-Reversion Traders* (Contrarians), and *Macro-Fundamental Anchors* (Curmudgeons).

\subsection{Main Contributions}
This paper establishes four key theoretical and empirical insights:
\begin{enumerate}
    \item \textbf{Impossibility of Herding in Isolated Markets}: We prove theoretically and empirically that market price signals evaluated in isolation ($p_{\text{peer}}=0$) cannot trigger herding or bistable potential wells ($k>0$ across all $M, \eta$). Independent cognitive errors cancel out under the Law of Large Numbers, proving that inter-agent sentiment transmission is a necessary condition for market regime shifts.
    \item \textbf{Empirical Verification of Critical Memory Depth ($M_c$)}: On real Nifty 50 Index returns, social peer transmission ($p_{\text{peer}}=1.0$) induces a pitchfork bifurcation between $M=9$ and $M=11$ days, matching the analytical mean-field prediction $M_c = \frac{\pi}{2(1-2\eta)^2} \approx 9.82$ days for a baseline cognitive noise rate $\eta = 0.30$.
    \item \textbf{Critical Slowing Down & Panic Relaxation ($T_{1/2}$)}: Simulating panic recovery from unanimous negative sentiment ($n_+(0)=0$), we observe classic *critical slowing down* for flat-weighted conformists near $M_c$, where relaxation times diverge ($T_{1/2} \to \infty$). Conversely, recency-weighted memory kernels ($b>1.0$) cap relaxation at $\sim 30$ days, acting as cognitive circuit breakers.
    \item \textbf{Cognitive Bias Overload & Feedback Inversion ($\eta > 0.50$)}: We demonstrate that cognitive noise rates $\eta > 0.50$ systematically invert the sentiment feedback operator, transforming the crowd into an effective contrarian network that quenches bubbles and enforces market stability.
\end{enumerate}

The remainder of this paper is structured as follows. Section 2 formulates the mathematical agent architecture and memory weight kernels. Section 3 derives the mean-field Fokker-Planck kinetic equation and Landau-Ginzburg drift expansion. Section 4 presents empirical potential landscapes reconstructed from the Nifty 50 Index. Section 5 evaluates relaxation dynamics, memory bias sweeps, and fundamental anchoring. Section 6 analyzes cognitive noise asymmetry, and Section 7 concludes with implications for systemic risk regulation.

\section{Agent Microarchitecture \& Behavioral Memory Kernels}

Consider an agent-based financial economy populated by $N$ market participants. At discrete trading day $t \in \mathbb{N}$, each agent $i \in \{1, \dots, N\}$ adopts a binary trading stance $s_i(t) \in \{-1, +1\}$, where $+1$ represents a bullish (long) posture and $-1$ represents a bearish (short) posture.

\subsection{Informational Queue & Cognitive Noise Operator}
Each agent maintains an internal First-In-First-Out (FIFO) cognitive queue $\mathbf{b}_i(t) = [b_i(t), b_i(t-1), \dots, b_i(t-M+1)]^T \in \{-1, +1\}^M$ of length $M$. On day $t$, agent $i$ receives a candidate sentiment bit $\tilde{q}_i(t)$ sampled via a Bernoulli mixture model:
\begin{equation}
    \tilde{q}_i(t) = 
    \begin{cases}
        s_{\text{peer}}(t-1) & \text{with probability } p_{\text{peer}} \\
        q_{\text{market}}(t-1) & \text{with probability } 1 - p_{\text{peer}}
    \end{cases}
\end{equation}
where $q_{\text{market}}(t-1) = \text{sign}(R_{t-1})$ denotes the 1-day lagged excess return direction of the asset, and $s_{\text{peer}}(t-1) \sim \text{Bernoulli}(n_+(t-1))$ represents a random peer opinion sampled from the crowd's order parameter:
\begin{equation}
    n_+(t) = \frac{1}{N} \sum_{j=1}^N \mathbb{I}(s_j(t) = +1)
\end{equation}
To model bounded rationality, emotional misinterpretation, and noisy signal processing, the candidate bit passes through a Binary Symmetric Channel (BSC) with misinterpretation rate $\eta \in [0, 0.5]$:
\begin{equation}
    b_i(t) = \mathcal{N}_\eta \left( \tilde{q}_i(t) \right) = 
    \begin{cases}
        -\tilde{q}_i(t) & \text{with probability } \eta \\
        +\tilde{q}_i(t) & \text{with probability } 1 - \eta
    \end{cases}
\end{equation}

\subsection{Behavioral Weight Kernels}
An agent evaluates its internal psychological sentiment score $\sigma_i(t)$ by computing a weighted inner product over its cognitive queue:
\begin{equation}
    \sigma_i(t) = \mathbf{w}_i^T \mathbf{b}_i(t) = \sum_{k=0}^{M-1} w_k \cdot b_i(t-k)
\end{equation}
where $\mathbf{w}_i$ is a normalized kernel satisfying $\sum_{k=0}^{M-1} |w_k| = 1$. The discrete voting mapping is given by $s_i(t) = \text{sign}(\sigma_i(t))$, defaulting to $+1$ if $\sigma_i(t) = 0$. We model five specific trader archetypes:

\begin{enumerate}
    \item \textbf{Herding Conformist (Pushover)}: Equal weight across memory depth $M$, representing unweighted crowd consensus:
    \begin{equation}
        w_k = \frac{1}{M}, \quad \forall k \in \{0, \dots, M-1\}
    \end{equation}
    \item \textbf{Recency-Biased Trader (Opportunist)}: Exponential weight decay parameter $b_{\text{opp}} > 1.0$, capturing availability heuristics and momentum chasing:
    \begin{equation}
        w_k = \frac{b_{\text{opp}}^k}{\sum_{j=0}^{M-1} b_{\text{opp}}^j}, \quad k \in \{0, \dots, M-1\}
    \end{equation}
    \item \textbf{Anchored History Trader (Traditionalist)}: Exponential decay parameter $0 < b_{\text{trad}} < 1.0$, modeling conservatism bias and historical anchoring:
    \begin{equation}
        w_k = \frac{b_{\text{trad}}^k}{\sum_{j=0}^{M-1} b_{\text{trad}}^j}, \quad k \in \{0, \dots, M-1\}
    \end{equation}
    \item \textbf{Mean-Reversion Trader (Contrarian)}: Inverts the internal memory majority sign, expecting price reversals:
    \begin{equation}
        s_i(t) = -\text{sign}\left( \sum_{k=0}^{M-1} \frac{1}{M} b_i(t-k) \right)
    \end{equation}
    \item \textbf{Macro-Fundamental Anchor (Curmudgeon)}: Disregards memory queues and peer chatter entirely, voting according to long-term macro trend derivatives (200-day moving average slope):
    \begin{equation}
        s_i(t) = C(t) = \text{sign}\left( \text{MA}_{200}(t-1) - \text{MA}_{200}(t-2) \right)
    \end{equation}
\end{enumerate}

\section{Mean-Field Theory \& Kinetic Bifurcation Mechanics}

To derive the macroscopic equation of motion for crowd sentiment, we apply a continuous-time mean-field approximation in the thermodynamic limit $N \to \infty$.

\subsection{Fokker-Planck Kinetic Derivation}
Let $x(t) = n_+(t) - 0.5 \in [-0.5, 0.5]$ represent crowd sentiment displacement from neutrality. In a purely peer-interacting crowd ($p_{\text{peer}}=1.0$), the probability that an agent receives a bullish bit $+1$ into its queue on day $t$ is:
\begin{equation}
    p_+ = (1-\eta) n_+(t-1) + \eta (1 - n_+(t-1)) = \frac{1}{2} + (1-2\eta) x(t-1)
\end{equation}
Over a memory queue of length $M$, the distribution of bullish entries follows a Binomial distribution $\mathcal{B}(M, p_+)$. Applying the Central Limit Theorem for $M \gg 1$, the continuous sentiment score $\sigma$ is Gaussian distributed:
\begin{equation}
    \sigma \sim \mathcal{N}\left( \mu_\sigma, \sigma_0^2 \right), \quad \mu_\sigma = M(2p_+ - 1) = 2M(1-2\eta)x, \quad \sigma_0^2 \approx M
\end{equation}
The probability that an individual agent votes bullish $+1$ is given by the error function integral:
\begin{equation}
    P(s_i = +1) = \frac{1}{\sqrt{2\pi \sigma_0^2}} \int_{0}^{\infty} e^{-\frac{(\sigma - \mu_\sigma)^2}{2\sigma_0^2}} d\sigma = \frac{1}{2} \left[ 1 + \text{erf}\left( \frac{\mu_\sigma}{\sqrt{2\sigma_0^2}} \right) \right]
\end{equation}
Substituting $\mu_\sigma$ and $\sigma_0$:
\begin{equation}
    n_+(t) = \frac{1}{2} + \frac{1}{2} \text{erf}\left( (1-2\eta) \sqrt{2M} \cdot x(t-1) \right)
\end{equation}
Subtracting $0.5$ yields the discrete order parameter update:
\begin{equation}
    x(t) = \frac{1}{2} \text{erf}\left( (1-2\eta) \sqrt{2M} \cdot x(t-1) \right)
\end{equation}

\subsection{Landau-Ginzburg Drift Expansion \& Critical Threshold $M_c$}
Expanding $\text{erf}(z) = \frac{2}{\sqrt{\pi}} \left( z - \frac{z^3}{3} + \mathcal{O}(z^5) \right)$ around $x=0$:
\begin{equation}
    \Delta x(t) = x(t) - x(t-1) \approx \left[ (1-2\eta) \sqrt{\frac{2M}{\pi}} - 1 \right] x - \frac{2}{3\sqrt{\pi}} \left( (1-2\eta)\sqrt{2M} \right)^3 x^3
\end{equation}
We define the Landau-Ginzburg drift equation:
\begin{equation}
    f(x) = \mathbb{E}[\Delta x \mid x] = -k x - c x^3
\end{equation}
where the linear curvature parameter $k$ and non-linear coefficient $c$ are:
\begin{equation}
    k = 1 - (1-2\eta) \sqrt{\frac{2M}{\pi}}, \quad c = \frac{4\sqrt{2}}{3\sqrt{\pi}} (1-2\eta)^3 M^{3/2}
\end{equation}
Integrating $f(x) = -\frac{dV}{dx}$ yields the effective psychological potential landscape:
\begin{equation}
    V(x) = \frac{1}{2} k x^2 + \frac{1}{4} c x^4
\end{equation}
A supercritical pitchfork bifurcation occurs precisely when $k=0$, defining the critical memory depth $M_c$:
\begin{equation}
    1 - (1-2\eta)\sqrt{\frac{2M_c}{\pi}} = 0 \implies M_c = \frac{\pi}{2(1-2\eta)^2}
\end{equation}
For our benchmark cognitive misinterpretation noise $\eta = 0.30$:
\begin{equation}
    M_c = \frac{\pi}{2(1 - 0.60)^2} = \frac{\pi}{2(0.16)} = \frac{\pi}{0.32} \approx 9.82 \text{ trading days}
\end{equation}
\begin{itemize}
    \item \textbf{Subcritical Regime ($M < M_c \implies k > 0$)}: $V(x)$ has a unique global minimum at $x=0$. Crowd sentiment is strictly mean-reverting and stable.
    \item \textbf{Supercritical Regime ($M > M_c \implies k < 0$)}: $x=0$ becomes an unstable local maximum. Two symmetric stable potential wells form at $x^* = \pm \sqrt{-k/c}$, driving spontaneous herding consensus.
\end{itemize}

\section{Empirical Potential Landscapes from Nifty 50 Data}

We test the analytical predictions using daily price return series from the Indian Nifty 50 Index (2,480 trading days). We run Monte Carlo simulations ($N=150$ agents) across memory depths $M \in [1, 21]$ and cognitive noise rates $\eta \in [0, 0.45]$.

\begin{figure}[htbp]
    \centering
    \includegraphics[width=0.85\linewidth]{Case_1a_-_Pure_Pushover_k_vs_M.png}
    \caption{Empirical curvature $k$ vs. cognitive memory depth $M$ calibrated on Nifty 50 daily returns ($\eta=0.30$). Isolated market processing ($p_{\text{peer}}=0$, blue) remains single-well ($k>0$) across all $M$. Peer-interacting crowds ($p_{\text{peer}}=1.0$, orange) cross $k=0$ between $M=9$ and $M=11$, confirming the analytical prediction $M_c = 9.82$.}
    \label{fig:k_vs_M}
\end{figure}

\subsection{Impossibility of Herding in Isolated Markets}
Figure \ref{fig:k_vs_M} compares the empirical curvature $k$ for isolated market processing ($p_{\text{peer}}=0$) versus peer-interacting crowds ($p_{\text{peer}}=1.0$). In isolated markets, $k$ remains strictly positive ($k > 0$) across all $M \in [1, 21]$. Table \ref{tab:market_only_grid} summarizes the empirical grid search across $\eta \times M$.

\begin{table}[htbp]
\centering
\caption{Empirical Curvature $k$ for Isolated Market Processing ($p_{\text{peer}}=0.0$). All parameter pairs exhibit $k > 0$ (single-well mean-reversion).}
\label{tab:market_only_grid}
\small
\begin{tabular}{@{}ccccccc@{}}
\toprule
$\eta \setminus M$ & $M=1$ & $M=5$ & $M=9$ & $M=13$ & $M=17$ & $M=21$ \\
\midrule
0.00 & +1.1913 & +0.3026 & +0.3026 & +0.3026 & +0.3026 & +0.3026 \\
0.10 & +1.1913 & +0.3026 & +0.3026 & +0.3026 & +0.3026 & +0.3026 \\
0.20 & +1.1913 & +0.3026 & +0.3026 & +0.3026 & +0.3026 & +0.3026 \\
0.30 & +1.1913 & +0.4106 & +0.4787 & +0.5072 & +0.4372 & +0.4718 \\
0.40 & +1.1913 & +0.4501 & +0.5123 & +0.5304 & +0.5102 & +0.5221 \\
\bottomrule
\end{tabular}
\end{table}

This empirical result proves that **external price signals alone cannot trigger financial herding**. Under independent sampling, cognitive noise is suppressed by the Law of Large Numbers ($N=150$), keeping market aggregate sentiment mean-reverting.

\begin{figure}[htbp]
    \centering
    \includegraphics[width=0.95\linewidth]{Case_1a_-_Pure_Pushover_potential.png}
    \caption{Reconstructed empirical potential landscapes $V(n_+)$ for Pushover crowds. Left: Isolated market processing ($p_{\text{peer}}=0$) maintains single-well restoring potentials. Right: Peer interaction ($p_{\text{peer}}=1.0$) develops bistable double-well energy landscapes for $M \ge 11$.}
    \label{fig:potential_landscapes}
\end{figure}

\subsection{Bistable Potential Landscapes under Peer Interaction}
Figure \ref{fig:potential_landscapes} displays the reconstructed potential profiles $V(n_+)$. In peer-interacting crowds, as cognitive memory expands past $M_c \approx 9.82$ days, the central state $n_+=0.5$ destabilizes, producing two symmetric equilibrium minima near $n_+ \approx 0.15$ (bearish panic) and $n_+ \approx 0.85$ (bullish euphoria).

\section{Panic Relaxation, Critical Slowing Down, \& Memory Bias}

To investigate crowd recovery from severe financial shocks, we conduct relaxation experiments starting from a unanimous bearish panic ($n_+(0)=0$, memory queues initialized to $-1$). We define the relaxation half-life $T_{1/2}$ as the trading days required for $n_+(t)$ to recover to neutrality ($0.50$).

\begin{figure}[htbp]
    \centering
    \includegraphics[width=0.85\linewidth]{thalf_vs_M.png}
    \caption{Relaxation half-life $T_{1/2}$ vs. cognitive memory depth $M$ calibrated on Nifty 50 data. Pushover crowds undergo critical slowing down as $M \to M_c$, failing to converge (DNF) for $M \ge 13$. Opportunists plateau rapidly ($T_{1/2} \approx 35$), while Contrarians flip instantly ($T_{1/2}=M$).}
    \label{fig:thalf_vs_M}
\end{figure}

\subsection{Critical Slowing Down}
Figure \ref{fig:thalf_vs_M} demonstrates classic *critical slowing down* in Herding Conformist crowds as memory approaches $M_c$:
\begin{equation}
    T_{1/2}(M=1)=3 \to T_{1/2}(M=5)=33 \to T_{1/2}(M=11)=1919 \text{ days}
\end{equation}
For $M \ge 13$ days, the crowd becomes permanently locked in panic consensus (Did Not Finish / DNF within the 2,480-day sample).

\subsection{Memory Bias Parameters \& Circuit Breaking}
We sweep the weight decay parameter $b$ for Opportunists ($b_{\text{opp}} \in [1.1, 2.0]$) and Traditionalists ($b_{\text{trad}} \in [0.5, 0.9]$).

\begin{figure}[htbp]
    \centering
    \includegraphics[width=0.85\linewidth]{thalf_heatmap_opportunist.png}
    \caption{Empirical heatmap of $T_{1/2}$ relaxation days for Opportunists across recency decay $b$ and memory depth $M$. High recency bias ($b=2.0$) caps panic recovery at $T_{1/2} \le 30$ days even for deep memory ($M=21$).}
    \label{fig:heatmap_opp}
\end{figure}

Figure \ref{fig:heatmap_opp} shows that strong recency weighting ($b_{\text{opp}}=2.0$) acts as an effective cognitive circuit breaker, preventing consensus lock-in and bounding recovery times ($T_{1/2} \le 30$ days) across all memory depths $M \in [1, 21]$.

\subsection{Macro Fundamental Anchoring (Curmudgeons)}
We evaluate the impact of introducing a minority fraction $\phi_C \in [0\%, 30\%]$ of Macro Fundamental Anchors into a Herding Conformist crowd.

\begin{figure}[htbp]
    \centering
    \includegraphics[width=0.85\linewidth]{thalf_heatmap_curmudgeon_nucleation.png}
    \caption{Curmudgeon minority nucleation effect on Pushover panic recovery $T_{1/2}$. Introducing $\ge 10\%$ fundamental anchors eliminates permanent lock-in traps ($\text{DNF}$) for $M \ge 13$, pinning sentiment recovery to $\sim 360$--$500$ trading days.}
    \label{fig:curmudgeon_nucleation}
\end{figure}

As shown in Figure \ref{fig:curmudgeon_nucleation}, introducing $\ge 10\%$ fundamental anchors eliminates the infinite lock-in trap ($\text{DNF}$), pinning market relaxation to $\sim 360$--$500$ trading days (corresponding to macro moving average trend reversals).

\section{Cognitive Noise Asymmetry Around $\eta = 0.50$}

The mean-field critical memory equation:
\begin{equation}
    M_c = \frac{\pi}{2(1-2\eta)^2}
\end{equation}
is mathematically symmetric around $\eta = 0.50$ due to the $(1-2\eta)^2$ term. However, our empirical simulations reveal a fundamental breakdown of this symmetry.

\subsection{Analytical Feedback Inversion}
The expected bit recorded by an agent sampling a peer is:
\begin{equation}
    \mathbb{E}[b_i(t)] = (2n_+ - 1)(1 - 2\eta)
\end{equation}
\begin{itemize}
    \item \textbf{Positive Feedback Regime ($\eta < 0.50 \implies 1-2\eta > 0$)}: High bullish crowd sentiment $n_+$ generates positive expected peer bits, amplifying sentiment displacement and driving **bistable herding ($k < 0$)**.
    \item \textbf{Negative Feedback Regime ($\eta > 0.50 \implies 1-2\eta < 0$)}: High bullish crowd sentiment $n_+$ generates negative expected peer bits due to noise inversion, creating an internal **restoring force ($k > 0$)**.
\end{itemize}

\begin{table}[htbp]
\centering
\caption{Empirical Curvature $k$ for Peer-Interacting Pushover Crowds ($p_{\text{peer}}=1.0$): $\eta=0.30$ vs. $\eta=0.70$.}
\label{tab:asymmetry}
\small
\begin{tabular}{@{}ccc@{}}
\toprule
Memory Depth $M$ & $\eta = 0.30$ (Positive Feedback) & $\eta = 0.70$ (Negative Feedback) \\
\midrule
1 & +0.6067 & +1.3849 \\
5 & +0.1304 & +0.4756 \\
9 & +0.0117 & +0.4401 \\
11 & \textbf{-0.0329 (BIFURCATION)} & +0.4160 \\
15 & \textbf{-0.0774 (BIFURCATION)} & +0.3794 \\
21 & \textbf{-0.0562 (BIFURCATION)} & +0.3336 \\
\bottomrule
\end{tabular}
\end{table}

As demonstrated in Table \ref{tab:asymmetry}, while $\eta = 0.30$ undergoes pitchfork bifurcation for $M \ge 11$, $\eta = 0.70$ remains strictly mean-reverting ($k > 0$) across all memory depths. High cognitive misinterpretation converts the crowd into an effective contrarian network, inherently suppressing financial bubbles.

\section{Discussion \& Financial Microstructure Implications}

Our empirical findings yield actionable implications for quantitative asset management and financial market regulation:

\begin{enumerate}
    \item \textbf{Monitoring Social Sentiment Contagion}: Regulators monitoring market risk must track social sentiment communication networks ($p_{\text{peer}}$) rather than price volatility alone, as sentiment contagion is the mandatory driver of market lock-in.
    \item \textbf{Cognitive Memory Horizons as Early Warning Indicators}: The threshold $M_c \approx 9.82$ trading days indicates that when market sentiment is driven by medium-term memory horizons ($\sim 2$ trading weeks), markets are highly susceptible to sudden regime shifts.
    \item \textbf{Role of Macro Trend Anchors}: Policies that incentivize long-term institutional fundamental investors ($\ge 10\%$ market share) prevent permanent liquidity traps and panic lock-in.
\end{enumerate}

\section{Conclusion}

This paper developed a behavioral memory-kernel agent-based model to analyze non-equilibrium phase transitions on the Indian Nifty 50 Index. We proved that isolated price processing cannot induce herding, whereas peer interaction triggers a pitchfork bifurcation at a critical cognitive memory depth $M_c \approx 9.82$ trading days. Furthermore, we quantified panic recovery dynamics, demonstrating critical slowing down in conformist crowds, circuit-breaking behavior in recency-biased momentum traders, and fundamental stabilization by macro anchors. Finally, we demonstrated that cognitive noise rates $\eta > 0.50$ invert sentiment feedback, providing a complete non-equilibrium framework for behavioral market dynamics.

\newpage
\bibliographystyle{IEEEtran}
\begin{thebibliography}{99}

\bibitem{fama1970efficient}
E.~F. Fama, ``Efficient capital markets: A review of theory and empirical work,'' \emph{The Journal of Finance}, vol. 25, no. 2, pp. 383--417, 1970.

\bibitem{shiller2000irrational}
R.~J. Shiller, \emph{Irrational Exuberance}.\hskip 1em plus 0.5em minus 0.4em\relax Princeton University Press, 2000.

\bibitem{lux1999scaling}
T.~Lux and M.~Marchesi, ``Scaling and criticality in a stochastic multi-agent model of financial markets,'' \emph{Nature}, vol. 397, no. 6719, pp. 498--500, 1999.

\bibitem{cont2001empirical}
R.~Cont, ``Empirical properties of asset returns: stylized facts and statistical issues,'' \emph{Quantitative Finance}, vol. 1, no. 2, pp. 223--236, 2001.

\bibitem{bouchaud2013more}
J.-P. Bouchaud, ``More is different in economics: the power of agent-based modeling,'' \emph{Nature Physics}, vol. 9, no. 9, pp. 519--520, 2013.

\bibitem{farmer2009economy}
J.~D. Farmer and D.~Foley, ``The economy needs agent-based modelling,'' \emph{Nature}, vol. 460, no. 7256, pp. 685--686, 2009.

\bibitem{kirman1993ants}
A.~Kirman, ``Ants, rational expectations, and recruitment in markets,'' \emph{Quarterly Journal of Economics}, vol. 108, no. 1, pp. 137--156, 1993.

\bibitem{bikhchandani1992theory}
S.~Bikhchandani, D.~Hirshleifer, and I.~Welch, ``A theory of fads, fashion, custom, and cultural change as informational cascades,'' \emph{Journal of Political Economy}, vol. 100, no. 5, pp. 992--1026, 1992.

\bibitem{Avery1998}
C.~Avery and P.~Zemsky, ``Multidimensional uncertainty and herd behavior in financial markets,'' \emph{American Economic Review}, vol. 88, no. 4, pp. 724--748, 1998.

\bibitem{kahneman1979prospect}
D.~Kahneman and A.~Tversky, ``Prospect theory: An analysis of decision under risk,'' \emph{Econometrica}, vol. 47, no. 2, pp. 263--291, 1979.

\bibitem{li2026informational}
X.~Li, Y.~Zhang, and M.~Rubenstein, ``Informational Memory Shapes Collective Behavior in Intelligent Swarms,'' \emph{Phys. Rev. Lett.}, vol. 136, p. 138302, 2026.

\end{thebibliography}

\end{document}
